\documentclass[10.5pt,a4paper]{article}

\usepackage{academic-cv}
\usepackage{fontawesome5}

\begin{document}

% ============================================================
% EN-TÊTE
% ============================================================

\cvname{Antoine LASSALLETTE}

\cvtitle{Enseignant supérieur --- Mathématiques appliquées \& Data Analytics}

\cvcontact{%
Lyon, France
\quad • \quad
\href{mailto:antoine_lassallette@msn.com}{antoine\_lassallette@msn.com}
\quad • \quad
\href{https://github.com/Lassallette}{github.com/Lassallette}
}

% ============================================================
% PROFIL
% ============================================================

\cvsection{Profil}

\cvprofile{%
Enseignant en mathématiques appliquées et Data Analytics, j'interviens dans
l'enseignement supérieur auprès d'étudiants en écoles de commerce, écoles
d'ingénieurs et formations professionnalisantes. J'enseigne en français comme
en anglais auprès de publics variés, du post-bac au Master.

Mon parcours associe plus de dix ans d'expérience pédagogique à une expérience
professionnelle de consultant en Data. Je privilégie une pédagogie progressive,
fondée sur la résolution de problèmes et les applications concrètes.
}

% ============================================================
% ENSEIGNEMENT SUPÉRIEUR ET POSTSECONDAIRE
% ============================================================

\cvsection{Expérience d'enseignement supérieur et postsecondaire}

% --- EM LYON ------------------------------------------------

\cventry
    {EM Lyon}
    {Intervenant en Mathématiques, Statistiques et Data Analytics}
    {2024--aujourd'hui}

\cvbullets{
    \item Enseignements en Global BBA et PGE :
    \textit{Data Analytics for Business}, mathématiques et statistiques
    appliquées, \textit{Maths et Stats 1}.

    \item \textit{Data Analytics for Business} : préparation et qualité des
    données, analyse exploratoire, bases de données, visualisation et
    formulation de recommandations fondées sur les données.

    \item Conception et adaptation de supports, exercices et évaluations ;
    accompagnement de projets d'analyse de données combinant SQL,
    modélisation relationnelle et visualisation.

    \item Enseignements dispensés en français et en anglais.
}

% --- ESME ---------------------------------------------------

\cventry
    {ESME}
    {Intervenant en Mathématiques et Analyse Numérique}
    {2024--2025}

\cvbullets{
    \item Cours et travaux dirigés de mathématiques en première année
    d'école d'ingénieurs : algèbre, analyse et outils mathématiques
    pour les sciences de l'ingénieur.

    \item Travaux dirigés d'analyse numérique, orientés vers la
    modélisation et les applications.

    \item Enseignements dispensés notamment en anglais auprès de
    groupes internationaux.
}

% --- LE WAGON -----------------------------------------------

\cventry
    {Le Wagon}
    {Intervenant en Data Science et Data Analytics}
    {2023--2024}

\cvbullets{
    \item Formation en mathématiques appliquées, statistiques, SQL,
    Python, Machine Learning et Data Analytics.

    \item Animation des enseignements et accompagnement des travaux
    pratiques, notebooks Python et projets de Data Science et
    Data Analytics.

    \item Intervention auprès d'apprenants adultes, principalement
    en reconversion professionnelle ou en complément de formation.
}

% --- SIMPLON ------------------------------------------------

\cventry
    {Simplon}
    {Formateur en Mathématiques et Intelligence Artificielle}
    {2024}

\cvbullets{
    \item Intervention auprès d'apprenants préparant le titre de
    Développeur en Intelligence Artificielle en alternance.

    \item Formation en statistiques appliquées à l'IA et au
    Machine Learning, notamment en régression et classification.

    \item Préparation et animation des enseignements, accompagnement
    des apprenants et évaluation des compétences.
}

% --- MADE IN ------------------------------------------------

\cventry
    {MADE iN}
    {Enseignant de Mathématiques --- BTS Design d'espace}
    {2018--2019}

\cvbullets{
    \item Enseignement des mathématiques générales auprès d'étudiants
    de BTS Design d'espace.

    \item Préparation des cours et exercices, accompagnement des
    étudiants et évaluation des apprentissages.
}

% ============================================================
% SAUT DE PAGE
% ============================================================

\newpage

% ============================================================
% ENSEIGNEMENT SECONDAIRE
% ============================================================

\cvsection{Enseignement secondaire}

\cventry
    {Enseignant de Mathématiques}
    {Collège et lycée}
    {2008--2021}

\cvdescription{%
Enseignement des mathématiques au collège et au lycée dans plusieurs
établissements en France et à l'étranger :
\textbf{Lycée Les Maristes} (Lyon, 2017--2021),
\textbf{Fénelon Sainte-Marie} (Paris, 2013--2017) et
\textbf{Saint Jean-Baptiste de La Salle} (Le Caire, Égypte, 2012--2013).
Enseignement auprès de classes de collège, lycée général et technologique,
incluant des classes à examen.
}

% ============================================================
% EXPÉRIENCE PROFESSIONNELLE
% ============================================================

\cvsection{Expérience professionnelle}

\cventry
    {BIPP}
    {Consultant Data}
    {2022--2023}

\cvbullets{
    \item Participation à des projets d'intégration et de migration de données
    au sein d'équipes techniques et métiers.

    \item Conception et suivi de flux ETL ; développement et maintenance
    de traitements de données avec Talend et SQL.

    \item Analyse des besoins utilisateurs et collaboration avec les équipes
    d'intégration et de restitution.

    \item Expérience des enjeux de la donnée en entreprise, aujourd'hui
    réinvestie dans mes enseignements en Data Analytics.
}

% ============================================================
% FORMATION
% ============================================================

\cvsection{Formation}

\cventry
    {Université Claude Bernard Lyon 1}
    {Licence Sciences et Technologies --- mention Mathématiques}
    {2005--2006}

\cventry
    {ISFA --- Université Claude Bernard Lyon 1}
    {Formation d'actuariat --- niveau Master 1}
    {2005--2008}

\cventry
    {CPGE MPSI / MP}
    {Formation scientifique supérieure}
    {2002--2005}

\cventry
    {CAFEP de Mathématiques}
    {Concours de recrutement de l'enseignement privé}
    {2008}

\cventry
    {Le Wagon}
    {Bootcamp Data Science}
    {2021}

% ============================================================
% DOMAINES D'ENSEIGNEMENT
% ============================================================

\cvsection{Domaines d'enseignement}

\cvline{Mathématiques appliquées :}
{Analyse, algèbre linéaire, analyse numérique, probabilités, statistiques,
équations différentielles, optimisation.}

\cvline{Data Analytics :}
{SQL, bases de données relationnelles, préparation et exploration des données,
Data Visualisation, Business Intelligence.}

\cvline{Data Science :}
{Python, statistiques appliquées, Machine Learning, Deep Learning.}

% ============================================================
% OUTILS
% ============================================================

\cvsection{Outils}

\cvline{Python \& Data Science :}
{NumPy, Pandas, Matplotlib, Scikit-learn, TensorFlow.}

\cvline{Data \& BI :}
{SQL, Tableau, Power BI, Talend, Excel, Hackolade.}

\cvline{Développement \& enseignement :}
{Git, GitHub, Jupyter, \LaTeX.}

\vspace{0.35cm}

\begin{center}
{\large
\faPython
\quad\textsf{Python}
\qquad\qquad
\faChartBar
\quad\textsf{Tableau}
\qquad\qquad
\textbf{\LaTeX}
}
\end{center}

\end{document}